\documentclass[10pt,a4paper]{amsart}
\usepackage[margin=19mm]{geometry}
\usepackage{amsmath,amssymb,mathtools}
\usepackage[scr=rsfs,cal=euler]{mathalfa}
\usepackage{xfrac}
\usepackage[unicode,hidelinks,bookmarks=false]{hyperref}
\newcommand{\ie}{i.e.\ }
\newcommand{\cf}{cf.\ }
\newcommand{\resp}{respectively\ }
\newcommand{\Subsection}{\subsection}
\providecommand{\nobreakdash}{\nobreak\hbox{-}}
\setlength{\emergencystretch}{2em}
\multlinegap0pt
\usepackage{bm}
\usepackage{bbm}
\usepackage{dsfont}
\usepackage{xspace}
\usepackage{cancel}
\usepackage{mathtools}
\usepackage{amsthm}
\makeatletter
\@ifundefined{lem}{\newtheorem{lem}{Lemma}}{}
\makeatother
\title[Real critical points of $T$-polynomials that are sums of squ]{Real critical points of $T$-polynomials that are sums of squared monomials and topology of $T$-hypersurfaces}
\author{Alo\"is Demory}
\date{Source: arXiv:2601.07751v1 (CC BY 4.0)}
\begin{document}
\maketitle

\noindent\emph{Source and licence.} Real critical points of $T$-polynomials that are sums of squared monomials and topology of $T$-hypersurfaces. Version arXiv:2601.07751v1 (CC BY 4.0), distributed under the Creative Commons Attribution 4.0 International licence, as recorded in the arXiv licence metadata for this exact version and shown on the abstract page https://arxiv.org/abs/2601.07751v1. Full rights notice: licenses/arxiv-2601.07751-CC-BY-4.0.txt.\par\smallskip

\noindent\textit{Editorial scope.} Counted unit: the Lem whose source label is \texttt{lemma\_points\_in\_cells} in the article (source0.tex lines 966--969, source label \texttt{lemma\_points\_in\_cells}), delivered together with the article's own complete original proof of it (source0.tex lines 971--994, one \texttt{proof} environment of 24 lines). No mathematics is written, repaired or re-derived: statement and proof are the article's own text line for line. Required context retained uncounted: none. Every definition, display and label of those lines is retained unchanged. Omitted from this package and not redistributed: 1--965, 970--970, 995--1078. Nothing inside a retained excerpt is omitted or altered: every retained body is delivered exactly as the article's lines are written, so no omission receipt of any kind accompanies this package. Citations: the retained excerpts carry no citation command at all, so no bibliography entry of the article is transcribed and no reference is redistributed. Reference adaptation: none was needed and none was made; every reference inside the delivered unit resolves inside it, so no unresolved reference or citation is produced. Printed slips kept literal: no printed word, symbol, hypothesis or number of the counted statement or of its proof was altered. Layout and class: the article's own class is \texttt{article} with packages including inputenc, fontenc, babel, xcolor, geometry, layout, setspace, charter, url, graphicx, tikz, amsmath, amssymb, mathrsfs; the wrapper is the lane's pinned \texttt{amsart} editorial wrapper at 10pt on A4, which restates the theorem-like environments the retained text needs and adds the article's own macro definitions that the retained text needs (none), with extra wrapper packages (bm, bbm, dsfont, xspace, cancel, mathtools). The delivered numbering is the unit's own. Assets: no retained span contains a figure environment, a graphics-inclusion command, a PDF-page inclusion, a table or a TikZ picture; no image file is copied into this package, the package declares no asset dependency and no image file is redistributed with it. Licence: version arXiv:2601.07751v1 of this article is distributed under the Creative Commons Attribution 4.0 International licence, as recorded in the arXiv licence metadata for this exact version and shown on the abstract page https://arxiv.org/abs/2601.07751v1. A full rights notice is delivered at licenses/arxiv-2601.07751-CC-BY-4.0.txt. Characters: every retained excerpt is pure ASCII, so the delivered engine, XeLaTeX with its default Latin Modern text font, reports no missing character.
\begin{lem}
    \label{lemma_points_in_cells}
    Every integer point lying in the interior of $\frac{m}{2}\Pi^n$ but not on a hyperplane with equation $\sum_{i=1}^n x_i = (n+1)k$ for a any integer $k$ lies in the interior of an $n$-dimensional cell of $\overline{\tau_m^n}$. Furthermore, the interior of every $n$-dimensional cell of $\overline{\tau_m^n}$ contains at most one integer point. 
\end{lem}

\begin{proof}
    Let $p = (p_1,..., p_n)$ be an integer point lying in the interior of $\frac{m}{2}\Pi^n$ but not on a hyperplane with equation $\sum_{i=1}^n x_i = (n+1)k$ for any integer $k$. Let $j$ be the integer in $\{1,...,n \}$ such that there exists a non-negative integer $k_0$ such that the point $p' := p - j(1,...,1)$ is on the hyperplane with equation $\sum_{i=1}^n x_i = (n+1)k_0$. Let $(k_1,...,k_n)$ be the unique solution of the following system of equations.
    \begin{equation}
    \begin{cases}
      2x_1 + \sum_{i \in \{2,...,n \}}x_i =  p_1 - j\\
      x_1 + 2x_2 + \sum_{i \in \{3,...,n \}}x_i = p_2 - j\\
      ...\\
      \sum_{i \in \{1,...,n-1\}}x_i + 2x_n = p_n-j
    \end{cases}
    \end{equation}
    The numbers $k_1,...,k_n$ are integers: the determinant of the matrix associated to the system is $n+1$ and the inverse of the matrix is $I_n - \frac{1}{n+1}J_n$, where $J_n$ is the matrix whose entries are all equal to $1$.
    Notice that, summing all the relations and dividing by $n+1$, we have $\sum_{i = 1}^n k_i = k_0$. Writing $p$ as $(2k_1 + k_2+...+k_n + j, k_1 + 2k_2 + k_3 + ... + k_n + j, ..., k_1 + ... + k_{n-1} + 2k_n + j)$, one can check that it belongs to the interior of the $n$-dimensional cell of $\overline{\tau_m^n}$ defined by the inequalities
    \begin{equation}
    \begin{cases}
        (n+1)(k_0 + j -1) < \sum_{i = 1}^n x_i < (n+1)(k_0 + j)\\
        -(n+1)(k_1-1) < -nx_1 + \sum_{i=2}^n x_i < -(n+1)k_1\\
        ...\\
        -(n+1)(k_n-1) < -nx_n + \sum_{i=1}^{n-1}n x_i < -(n+1)k_n
    \end{cases} \, .
    \end{equation}
    Indeed, the first line becomes $(n+1)(j-1) < nj < (n+1)j$ after substracting $(n+1)k_0 = \sum_{i=1}^n (n+1)k_i$, while the others all give $-(n+1) < -j < 0$.

    The second part of the lemma is then straightforward.
\end{proof}

\end{document}
